% 第二章 相关技术

\chapter{相关技术}

\section{桌面端与链下服务技术基础}

FreeFlow 选用 Electron 与 React 来构建客户端，并使用 Express 和 Prisma 来搭建链下服务。项目使用统一的Node 环境 MonoRepo 策略管理项目以提升开发效率。

Electron 提供了跨平台桌面应用的运行环境，同时还依靠主进程与渲染进程的分工来承载系统功能\cite{electronDocs,electronIPC}。在本系统中，主进程会负责窗口管理、本地数据库访问、下载与缓存控制、音频代理等本地能力方面的管理工作；渲染进程则负责创作者工作台、资源检索、详情展示、评论与播放控制等UI的实现，网页部分采用 React\cite{reactDocs} 构建。

Node.js 适宜处理上传、流式转发以及链上交互代理等 I/O 密集型任务，而 Express 则在项目中用于认证、资源、购买、评论以及异常处理等业务接口\cite{nodeDocs,expressDocs}实现。系统运用了 PostgreSQL 来保存用户、钱包身份、发行记录、上传记录、购买记录与评论数据，并且凭借 Prisma 来完成模式定义、关系映射以及迁移管理工作\cite{prismaDocs}。

\section{SIWE身份认证机制}

本文中的Web2.5服务器基于 SIWE（Sign-In with Ethereum）实现，符合 EIP-4361 定义的消息格式\cite{eip4361}。该机制凭借钱包签名来证明用户对链上地址的控制权，能够把系统身份与链上账户建立起统一的关联关系，方便后续的购买、发布以及访问校验等流程复用同一个身份基础。

在具体实现当中，后端会先生成一个一次性的随机数 nonce，并且构造出包含域名、地址、URI、链 ID、签发时间等字段的标准化消息；客户端调起钱包完成签名之后，再把签名结果和原始消息一起提交给后端；后端随后会对消息格式、nonce 有效性、链 ID 以及签名恢复地址进行校验\cite{eip4361}。一旦验证通过，系统便会绑定用户钱包地址并关联后续业务。

相较于传统账户密码登录，SIWE 使用证明钱包控制权的机制来避免中心化身份，同时更符合本系统的应用场景。

\section{IPFS内容寻址存储与混合架构}

链上音乐系统牵涉到媒体资源的存储、可验证的引用以及高频的业务协同。正是鉴于这一特性，FreeFlow 把区块链、IPFS 以及链下服务组合成了混合架构，以便兼顾可信记录、资源可定位性和系统运行效率。

IPFS 运用了内容寻址机制，用 CID 来标识文件内容。因此在FreeFlow中适合存储音频文件、封面图以及 metadata 等静态对象\cite{ipfsDocs}。本系统采用了 Pinata 来完成上传与持久化的管理工作，在创作者发布的流程当中会保存音频、封面和 metadata 对象，并且获得对应的 CID 或网关地址\cite{pinataDocs}。考虑到 IPFS 本身主要负责寻址与分发方面的工作，资源的长期可用性还是要依赖于持续的保存机制，系统使用 Pinata 的持久化服务来降低节点维护以及 pinning 管理方面的复杂度\cite{ipfsPinning}。

在系统层面，作品发布结果、访问凭证、交易记录以及收益规则等关键状态会被保存到链上；而评论、资源索引、登录会话、上传状态与检索支撑数据则保存于链下。已经有研究指出，链上、链下以及混合模式在成本、吞吐、可扩展性还有可信性之间存在着明确的权衡关系，面向实际应用的区块链系统通常需要把链上记录与链下处理结合起来开展协同设计\cite{Wei2022BDMSSurvey,Six2022DAppPatterns,Liu2020GasExceptions}。对音乐场景，音轨等资源文件体积较大，评论与数据检索信息更新频繁。若全量上链，成本与时间不可控；把高可信状态和高频服务进行分层组织，更加契合系统当前的规模与实现目标。

\section{智能合约与ERC-1155访问凭证}

链上业务逻辑选用了 Solidity 来进行实现\cite{solidityDocs}。系统依据作品发布、访问授权以及收益分配这几个方面来设计合约结构，其中的访问控制以 ERC-1155 多代币标准作为基础\cite{eip1155}合约接口。

ERC-1155 允许在同一份合约里面管理多种代币类型，可把每首作品映射成独立的\texttt{tokenId}访问凭证。在 FreeFlow 中，用户完成购买以后会获得对应的\texttt{tokenId}，系统依据\texttt{tokenId}进行访问状态校验。

系统的合约主要包括了 \texttt{MusicAccess1155}、\texttt{PlatformHub} 以及 \texttt{RoyaltySplitter}。\texttt{MusicAccess1155} 负责访问凭证的创建与校验工作；\texttt{PlatformHub} 负责作品发布、销售参数维护、购买处理以及访问判断方面的工作；\texttt{RoyaltySplitter} 则负责收益的分账与提取操作。为了提升标准实现的可靠性，开发过程中复用了 OpenZeppelin 所提供的一系列合约组件与安全工具库\cite{openzeppelinContracts}。

\section{音频代理、Range请求与本地缓存}

在线播放场景对于随机访问、进度跳转以及重复播放效率方面有着较高的要求。为了提高桌面端的播放稳定性，系统在本地集成了音频代理服务，同时把本地文件、缓存文件以及远程音频统一都纳入到同一个访问入口里。

播放器通过本地代理接口来请求音频数据。代理层收到请求后，根据资源缓存状态判断直接读取本地文件、返回缓存内容，又或者是向远程地址发起转发请求。对于客户端界面中拖动进度条、断点续播等用户行为，系统依赖 HTTP Range 机制完成读取工作。依据 HTTP 语义规范，客户端可以凭\texttt{Range}请求头来指定所要获取音轨的播放区间，服务端会返回\texttt{206 Partial Content}响应\cite{rfc9110}。这种机制在FreeFlow中有助于降低重新加载完整音频文件所带来的时间方面的开销。

在实现上，Node.js 所提供的 HTTP 与流式处理方面的能力适合 FreeFlow 中的音频转发和缓存写入这类场景\cite{nodeHttp,nodeStream}。系统针对普通的播放请求会采用边传输边缓存的策略，而对于范围请求则会优先顺延传递上游响应，在保证跳转实时性的同时，也兼顾到后续的复用效率。